% === Begin Label Data ===
% ID: 500
% 难度星级: 
% 标签: 
% 备注: 
% 组卷引用次数: 0
% === End  Label Data ===

\begin{problem}{2020}{G}{新高考I卷}{6}{函数}
基本再生数 $ R_0 $ 与世代间隔 $ T $ 是新冠肺炎流行病学基本参数．基本再生数指一个感染者传染的平均人数，世代间隔指相邻两代间传染所需的平均时间．在新冠肺炎疫情初始阶段，可以用指数模型：$ I(t)=e^{rt} $ 描述累计感染病例数 $ I(t) $ 随时间 $ t $（单位：天）的变化规律，指数增长系数 $ r $ 与 $ R_0 $，$ T $ 近似满足 $ R_0=1+rT $．有学者基于已有数据估计出 $ R_0=3.28 $，$ T=6 $．据此，在新冠肺炎疫情初始阶段，累计感染病例数增加 $ 1 $ 倍需要的时间约为（$ \ln 2\approx 0.69 $）
\begin{choices}
\choice{{${1.2}$ 天}}
\choice{{${1.8}$ 天}}
\choice{{${2.5}$ 天}}
\choice{{${3.5}$ 天}}
\end{choices}
\end{problem}

\begin{answer}

\end{answer}

\begin{solutions}

\end{solutions}